\chapter{Discussion and Conclusions}\label{ch:conclusions}

\section{Findings}\label{sec:conclusions-findings}\label{sec:conclusions-amp}

\paragraph{Derived exactly.} For the stationary AR(1) chain under
coordinatewise OU corruption, the joint score is $\score(x,t) = -\Qt x$
with $\Qt = (e^{-2t}\Sigz + \Deltat I)^{-1}$, equal by the
score--posterior identity to $(e^{-t}\E[a|x] - x)/\Deltat$ (RQ1). Its
coupling structure has a closed-form lifecycle: exactly tridiagonal at
$t = 0$, filling one diagonal per order in $t$ by the band-fill law
\eqref{eq:g-bandfill}, and returning to the identity at rate $e^{-2t}$.
In the bulk, the single quantity $q(\alpha, t) \in (0, |\alpha|]$ governs
both the posterior correlation length and the geometric decay of evidence
influence, with limits $q \to 0$ as $t \to 0$ and $q \to |\alpha|$ as
$t \to \infty$. The posterior of the noisy chain lives on a tree-shaped
factor graph, on which one forward and one backward BP sweep compute all
smoothing marginals; on the Gaussian chain the messages close
algebraically in the Gaussian family, and the recursion is the Kalman
filter followed by the RTS smoother (RQ2). The
error of a radius-$r$ local score estimator decays exactly as $q^r$
(RQ4). On the same model, the per-node AMP/TAP closure returns the exact
score (the mean fixed point is closure-independent) but a different
variance, with existence boundary, breakdown time $t_c(\alpha)$, and
critical coupling $\alpha_c = \sqrt{2} - 1$ in closed form
(Appendix~\ref{app:amp}); all of these formulas were re-verified
independently.

\paragraph{Measured numerically.} The two Gaussian score routes agree to
$10^{-15}$--$10^{-14}$ across the tested $(K, \alpha, t)$ grid, with an
independently written implementation reproducing the agreement. On
Laplace chains, Gaussian-projected BP measured against validated grid
quadrature has median relative score error $0.39$ at $t = 0.02$,
decaying below $10^{-3}$ for $t \gtrsim 0.9$; the error is amplified at
small $t$ by the exact factor $e^{-t}/\Deltat$, is worst for bimodal
innovation laws, and decreases with stronger chain correlation, while
the score direction remains accurate (median cosine $\geq 0.92$
throughout) (RQ3). At equal locality budget, truncating the inference is
more accurate than truncating the score matrix at small $t$ ($12.3\%$
against $21.0\%$ at $t = 0.05$).

\paragraph{The Gaussian-to-Laplace comparison.} The structural lesson is
where exactness is lost: the tree, the two-sweep schedule, and the
score--posterior identity survive any innovation law; what breaks is the
representability of messages in a finite-dimensional family. The
practical lesson is when it matters: the Gaussian closure is essentially
free at moderate and large diffusion times, and expensive precisely in
the small-$t$ regime where the score magnitude is amplified, though even
there the score direction is preserved much better than its magnitude.

\paragraph{Learning the innovation law.} Chapters~\ref{ch:em-method}
and~\ref{ch:em-results} relax the assumption, held throughout the rest of
this thesis, that the innovation law is known. An EM procedure built
directly on the exact chain inference of Chapter~\ref{ch:gaussian} estimates
a mixture innovation from noised data alone, with an exact E-step for any
innovation law and a generalised M-step; the resulting estimator reduces the
data needed for accurate denoising by $9$--$14\times$ against an unstructured
neural network and $2$--$4\times$ against a locality-matched one. That
pointwise advantage does not transfer to generation: the sharpest evidence
is a single trained network becoming simultaneously more accurate pointwise
and less faithful generatively as its own training data grows, so pointwise
accuracy does not determine generative fidelity even within one estimator.
Mixture capacity saturates, on held-out evidence and on paired pointwise
error, at roughly eight components. The chain-Markov assumption itself
tolerates a violation it can absorb into its own fitted parameters (a
shared global latent, up to the point where it accounts for half the
marginal variance) and fails, well before generic intuition about
misspecification would suggest, on one it structurally cannot represent
(long-range coupling, by a coupling strength of roughly one tenth).

\section{What the model suggests but does not establish}

Within the Gaussian benchmark, local score computations are near-optimal
away from intermediate diffusion times, the required receptive field
scales as $\xi \log(1/\varepsilon)$ with $\xi = 1/\log(1/q)$, and a
local algorithm is preferable to a banded linear readout. These are
properties of exact inference in one controlled model. Whether trained
score networks on real sequence data inherit them is untested here and
would require learning experiments.

\section{Limitations}

(i) The exact solutions concern one-dimensional state per frame and
linear AR(1) dynamics, extended to a two-dimensional nonlinear state only
in Chapter~\ref{ch:ring}. (ii) The non-Gaussian results are numerical: grid
BP is a reference only within its validated resolution regime
($\sqrt{2t} \gtrsim 3\,\dd x$), and the Laplace closure error is a
measured curve, not a formula. (iii) The locality laws are proved for
the Gaussian bulk; their non-Gaussian counterparts are only structurally
delimited. (iv) An earlier conjecture on low-rank structure of the Laplace
$K = 2$ posterior Hessian failed outside a narrow regime and is not
asserted. (v) The EM-learning results of Chapters~\ref{ch:em-method}
and~\ref{ch:em-results} use a fixed, generous but ultimately insufficient
iteration budget for several of the reported sweeps; any conclusion in
those chapters resting on fitted or generated \emph{shape} at that budget
is flagged explicitly as confounded with convergence rate rather than
resolved. (vi) The misspecification results of
Section~\ref{sec:em-nonmarkov} probe exactly two violations of the Markov
assumption; they establish that the two are not interchangeable in effect,
not a general theory of how much non-Markov structure a chain estimator
tolerates. (vii) The image and video extensions of the EM/BP estimator,
explored in the underlying research package, are outside the scope of
this thesis and are not reported here.

\section{Future work}

Two of the three directions identified earlier in this project have since
been carried out and are reported in Chapters~\ref{ch:em-method}
and~\ref{ch:em-results}: mixture messages, in the form of a fully estimated
$C$-component innovation kernel rather than a fixed two-component closure;
and reverse dynamics, in the form of the generative comparison of
Section~\ref{sec:em-generative}, which integrates the reverse SDE under the
estimated EM/BP score and under two trained neural scores. Discrete chains
remain untouched by this thesis: with a finite alphabet the BP messages are
finite vectors and inference is exact with no closure at all, which would
give a second solvable benchmark alongside the Gaussian one.

What the new chapters measured opens further directions of their own.
Re-running the capacity sweep of Section~\ref{sec:em-capacity} to a
convergence criterion rather than a fixed iteration count, now that
Section~\ref{sec:em-convergence} has measured what that criterion should be,
would settle the one result in this thesis that is explicitly flagged as
confounded. The misspecification study of Section~\ref{sec:em-nonmarkov}
probes two violations at a time; a paired design across violation strengths,
and further violation types beyond a global latent and long-range coupling,
would generalise the finding that the two mechanisms are not
interchangeable. And the reverse-dynamics comparison itself used a fixed,
small chain length and one innovation family; whether the pointwise/generative
dissociation of Section~\ref{sec:em-generative} scales with sequence length,
and whether it appears under other non-Gaussian innovation laws, remain
open.
